% =====================================================================
%  Capitolo 3 — Fondamenti teorici
% =====================================================================
\chapter{Fondamenti teorici}
\label{cap:teoria}

Questo capitolo sviluppa il quadro teorico dei meccanismi studiati. Si parte dalla formulazione
del problema e dalla funzione di perdita, poi si passa all'architettura encoder-decoder e ai
meccanismi di attenzione, fino alla derivazione dell'attenzione lineare usata da \maresunet{}.
Seguono la quantizzazione vettoriale e l'auto-distillazione, la fusione dell'informazione
altimetrica, le metriche e l'inferenza su tile intere. Per ogni meccanismo si mettono in
evidenza le proprietà che motivano le domande di ricerca e le scelte del
\cref{cap:metodologia}.

% ---------------------------------------------------------------------
\section{Formulazione del problema}
\label{sec:formulazione}

Un'immagine è un tensore $\vect{x}\in\R^{H\times W\times C_{\mathrm{in}}}$, con $C_{\mathrm{in}}=3$
(RGB o IR-R-G) oppure $C_{\mathrm{in}}=4$ se si aggiunge l'nDSM. L'insieme delle posizioni è
$\Omega=\{1,\dots,H\}\times\{1,\dots,W\}$. L'etichetta è una mappa
$y:\Omega\to\{1,\dots,K\}\cup\{\varnothing\}$, dove $K=6$ è il numero di classi e
$\varnothing$ (codificato con il valore 255) indica un pixel senza etichetta valida. Le
posizioni etichettate sono $\Omega_v=\{u\in\Omega : y_u\neq\varnothing\}$.

Un modello di segmentazione $f_\theta$ associa a ogni posizione un vettore di logit
$\vect{z}_u=f_\theta(\vect{x})_u\in\R^K$ e quindi la distribuzione
$p_\theta(k\mid\vect{x},u)=\softmax(\vect{z}_u)_k$. L'addestramento minimizza l'entropia
incrociata media sulle posizioni etichettate:
\begin{equation}
  \mathcal{L}_{\mathrm{CE}}(\theta)
  = -\frac{1}{\abs{\Omega_v}}\sum_{u\in\Omega_v}\log p_\theta(y_u\mid\vect{x},u).
  \label{eq:ce-mascherata}
\end{equation}
La predizione è $\hat y_u=\argmax_k \vect{z}_{u,k}$. La scelta di \emph{quali} pixel entrano in
$\Omega_v$ non è un dettaglio implementativo, come mostra il risultato seguente.

\begin{proposizione}[Classe esclusa dalla perdita]
\label{prop:classe-esclusa}
Si escluda dalla perdita~\eqref{eq:ce-mascherata} una classe reale $c$, cioè si sostituisca
$\Omega_v$ con $\Omega_v'=\{u\in\Omega_v : y_u\neq c\}$. Allora, per ogni posizione
$u\in\Omega_v'$,
\[
  \frac{\partial\mathcal{L}}{\partial z_{u,c}}=\frac{p_\theta(c\mid\vect{x},u)}{\abs{\Omega_v'}}\;\geq\;0,
\]
mentre i pixel di classe $c$ non contribuiscono alla perdita. Ogni passo di discesa del gradiente
riduce quindi il logit della classe $c$ e non lo aumenta mai.
\end{proposizione}
\begin{proof}
Per una singola posizione $\ell_u=-\log\softmax(\vect{z}_u)_{y_u}$ si ha
$\partial\ell_u/\partial z_{u,k}=p_\theta(k\mid\vect{x},u)-\Ind[k=y_u]$. Per $u\in\Omega_v'$ vale
$y_u\neq c$, quindi la derivata rispetto a $z_{u,c}$ è $p_\theta(c\mid\vect{x},u)\ge0$. Le
posizioni con $y_u=c$ non compaiono nella somma.
\end{proof}

Nel codice del 2022 l'indice ignorato era 0, che corrispondeva proprio al clutter. Per la
\cref{prop:classe-esclusa} il modello riceveva solo una pressione a \emph{non} prevedere il
clutter. Poteva ancora prevederlo dove l'ingresso era diverso da tutti i pixel etichettati,
perché lì la pressione era solo indiretta: questo è compatibile con l'IoU basso ma non nullo
riportato allora ($\approx0{,}15$). In questo lavoro si ignora soltanto il valore $\varnothing$.

% ---------------------------------------------------------------------
\section{Encoder-decoder residuali}
\label{sec:encoder-decoder}

\paragraph{Blocchi residui.}
Un blocco residuo~\cite{he2016resnet} calcola $\vect{x}_{l+1}=\vect{x}_l+\mathcal{F}(\vect{x}_l;W_l)$.
Nella forma semplificata con mappa identità, srotolando la ricorrenza si ottiene
$\vect{x}_L=\vect{x}_l+\sum_{i=l}^{L-1}\mathcal{F}(\vect{x}_i;W_i)$ e quindi
\begin{equation}
  \frac{\partial\mathcal{L}}{\partial\vect{x}_l}
  =\frac{\partial\mathcal{L}}{\partial\vect{x}_L}
   \Bigl(I+\frac{\partial}{\partial\vect{x}_l}\sum_{i=l}^{L-1}\mathcal{F}(\vect{x}_i;W_i)\Bigr).
\end{equation}
Il termine identità fa arrivare il gradiente agli strati iniziali senza attenuazione
moltiplicativa. Per questo encoder profondi pre-addestrati si possono rifinire in modo stabile.

\paragraph{Gerarchia di scale.}
Gli encoder ResNet producono quattro mappe $\vect{E}_1,\dots,\vect{E}_4$ a stride $4,8,16,32$
(\cref{tab:stadi-resnet}). Le mappe superficiali sono ad alta risoluzione e semanticamente
povere, quelle profonde sono povere di dettaglio spaziale. Il numero di posizioni
$N_l=(H/s_l)(W/s_l)$ diminuisce di un fattore 4 a ogni stadio. È questa quantità a governare il
costo dell'attenzione (\cref{sec:attenzione-dot}).

\begin{table}[t]
  \centering
  \small
  \caption[Stadi degli encoder ResNet per un ritaglio di 512 pixel]{Stadi degli encoder ResNet
    per un ritaglio di $512\times512$ pixel: stride, canali, numero di posizioni $N_l$ e memoria
    della sola matrice di attenzione $N_l\times N_l$ in float32, per immagine e per testa.}
  \label{tab:stadi-resnet}
  \begin{tabular}{|c|c|c|c|c|r|}
    \hline
    \textbf{Stadio} & \textbf{Stride} & \makecell{\textbf{Canali}\\\textbf{R18/R34}} & \makecell{\textbf{Canali}\\\textbf{R50}} &
    \textbf{Posizioni $N_l$} & \makecell{\textbf{Memoria}\\\textbf{$N_l^2$}} \\
    \hline
    $\vect{E}_1$ & 4  & 64  & 256  & $128^2=\num{16384}$ & \SI{1}{GiB} \\
    $\vect{E}_2$ & 8  & 128 & 512  & $64^2=\num{4096}$   & \SI{64}{MiB} \\
    $\vect{E}_3$ & 16 & 256 & 1024 & $32^2=\num{1024}$   & \SI{4}{MiB} \\
    $\vect{E}_4$ & 32 & 512 & 2048 & $16^2=256$          & \SI{256}{KiB} \\
    \hline
  \end{tabular}
\end{table}

\paragraph{Decoder con connessioni di salto.}
Uno schema U-Net~\cite{ronneberger2015unet} ricostruisce la risoluzione combinando, a ogni scala,
l'uscita sovracampionata del livello inferiore con una trasformazione $\mathcal{S}_l$ della
mappa dell'encoder alla stessa scala:
\begin{equation}
  \vect{D}_l=\mathrm{Dec}_{l+1}(\vect{D}_{l+1})+\mathcal{S}_l(\vect{E}_l).
  \label{eq:decoder}
\end{equation}
Nella U-Net \enquote{semplice} $\mathcal{S}_l$ è l'identità. In \maresunet{} è un blocco di
attenzione (\cref{sec:maresunet}). È esattamente questa differenza a essere misurata dalla
\RQ{1}.

% ---------------------------------------------------------------------
\section{Meccanismi di attenzione}
\label{sec:attenzione}

\subsection{Attenzione dot-product}
\label{sec:attenzione-dot}

Sia $\vect{X}\in\R^{N\times C}$ una mappa di caratteristiche con le $N=hw$ posizioni disposte in
righe. L'attenzione dot-product~\cite{vaswani2017attention} proietta $\vect{X}$ in query, key e
value, $Q=\vect{X}W_q\in\R^{N\times d_k}$, $K=\vect{X}W_k\in\R^{N\times d_k}$ e
$V=\vect{X}W_v\in\R^{N\times C}$, e calcola
\begin{equation}
  D(Q,K,V)=\softmax_{\mathrm{riga}}\!\left(\frac{QK\T}{\sqrt{d_k}}\right)V .
  \label{eq:dot-product}
\end{equation}
L'uscita in posizione $i$ è una media pesata dei value di \emph{tutte} le posizioni, con pesi
dati dalla somiglianza tra la query $\vect{q}_i$ e le key $\vect{k}_j$. La matrice $QK\T$ ha
$N^2$ elementi: calcolo e memoria sono $\bigO{N^2 d_k}$ e $\bigO{N^2}$. Sul primo stadio di un
ritaglio di $512\times512$ pixel ($N=\num{16384}$) la sola matrice occupa \SI{1}{GiB} in
float32 per immagine e per testa (\cref{tab:stadi-resnet}). Con un batch di 8 immagini si
superano gli \SI{8}{\giga\byte} della GPU usata in questo lavoro prima ancora di contare le
attivazioni del resto della rete (\cref{fig:complessita}).

\begin{figure}[t]
  \centering
  % Memoria richiesta dall'attenzione sul primo livello dell'encoder (stride 4, C = 64,
% d_k = C/8) in funzione del lato del ritaglio, per una singola immagine in float32.
\begin{tikzpicture}
  \begin{axis}[
      width=0.9\linewidth, height=6.2cm,
      xlabel={lato del ritaglio $s$ (pixel)},
      ylabel={memoria (GiB, scala log.)},
      xmin=128, xmax=1024, ymode=log, ymin=1e-5, ymax=60,
      xtick={128,256,384,512,640,768,896,1024},
      grid=major, grid style={SapGrigioChiaro},
      legend pos=north west, legend cell align=left,
      legend style={font=\footnotesize, draw=SapGrigioTenue},
      tick label style={font=\footnotesize}, label style={font=\small},
      samples=60, domain=128:1024,
    ]
    % N = (s/4)^2 ; dot-product: N^2 * 4 B ; LAM: (N (C + d_k) + d_k C) * 4 B
    \addplot[SapRosso, very thick] {((x/4)^2)^2*4/2^30};
    \addlegendentry{dot-product: matrice $N\times N$}
    \addplot[SapGrigio, very thick, dashed] {(((x/4)^2)*(64+8) + 8*64)*4/2^30};
    \addlegendentry{LAM: stato $d_k\times C$ e uscite $N\times(C+d_k)$}
    \addplot[black, thin, densely dotted, domain=128:1024] {8};
    \node[font=\scriptsize, anchor=south east] at (axis cs:1020,8) {\SI{8}{\giga\byte} (RTX 4060)};
    \addplot[only marks, mark=*, mark size=1.6pt, SapRossoScuro] coordinates {(512,1)};
    \node[font=\scriptsize, anchor=west, text=SapRossoScuro] at (axis cs:525,0.6) {\SI{1}{GiB} per immagine a $s=512$};
  \end{axis}
\end{tikzpicture}

  \caption[Memoria dell'attenzione dot-product e lineare]{Memoria richiesta dall'attenzione sul
    primo stadio dell'encoder (stride 4, $C=64$, $d_k=C/8$) al variare del lato del ritaglio,
    per una singola immagine in float32. La matrice $N\times N$ cresce con la quarta potenza
    del lato. Lo stato dell'attenzione lineare cresce con il quadrato, cioè linearmente in
    $N$.}
  \label{fig:complessita}
\end{figure}

\subsection{Attenzione come media a nucleo}
\label{sec:attenzione-nucleo}

La riga $i$ della~\eqref{eq:dot-product} si può scrivere come
\begin{equation}
  D_i=\frac{\sum_{j=1}^{N}\mathrm{sim}(\vect{q}_i,\vect{k}_j)\,\vect{v}_j}
           {\sum_{j=1}^{N}\mathrm{sim}(\vect{q}_i,\vect{k}_j)},
  \qquad \mathrm{sim}(\vect{q},\vect{k})=\exp\!\bigl(\vect{q}\T\vect{k}/\sqrt{d_k}\bigr),
  \label{eq:attenzione-nucleo}
\end{equation}
cioè come una media a nucleo dei value. La~\eqref{eq:attenzione-nucleo} definisce
un'attenzione valida per qualunque funzione di somiglianza non negativa. Se
$\mathrm{sim}(\vect{q},\vect{k})=\phi(\vect{q})\T\psi(\vect{k})$ per due mappe
$\phi,\psi:\R^{d_k}\to\R^{m}$, la proprietà associativa del prodotto
matriciale~\cite{katharopoulos2020linear} consente di raccogliere le somme sulle key, che non
dipendono da $i$:
\begin{equation}
  D_i=\frac{\phi(\vect{q}_i)\T\Bigl(\sum_{j}\psi(\vect{k}_j)\,\vect{v}_j\T\Bigr)}
           {\phi(\vect{q}_i)\T\sum_{j}\psi(\vect{k}_j)} .
  \label{eq:attenzione-lineare-generale}
\end{equation}
La matrice $\sum_j\psi(\vect{k}_j)\vect{v}_j\T\in\R^{m\times C}$ e il vettore
$\sum_j\psi(\vect{k}_j)\in\R^{m}$ si calcolano una sola volta in $\bigO{NmC}$. Ogni uscita
costa poi $\bigO{mC}$: il costo complessivo è lineare in $N$ (\cref{fig:associativita}).

\begin{figure}[t]
  \centering
  \resizebox{0.92\linewidth}{!}{% Ordine dei prodotti: (Q K^T) V costa O(N^2 d), Q (K^T V) costa O(N d C).
\begin{tikzpicture}[x=1mm, y=1mm, font=\footnotesize]
  \tikzset{mat/.style={draw=SapRossoScuro, line width=0.5pt}}
  % --- sinistra: (Q K^T) V -------------------------------------------------
  \begin{scope}
    \draw[mat, fill=SapRossoChiaro] (0,0) rectangle (4,24);      \node at (2,-3) {$Q$};
    \node at (6,12) {$\cdot$};
    \draw[mat, fill=SapRossoChiaro] (8,20) rectangle (32,24);    \node at (20,17) {$K\T$};
    \node at (36,12) {$=$};
    \draw[mat, fill=SapRossoTenue] (40,0) rectangle (64,24);     \node at (52,12) {$N\times N$};
    \node at (68,12) {$\cdot$};
    \draw[mat, fill=SapGrigioChiaro, draw=SapGrigio] (72,0) rectangle (78,24); \node at (75,-3) {$V$};
    \node[text=SapRossoScuro, font=\footnotesize\bfseries] at (39,31) {$(QK\T)\,V$: memoria e calcolo $\mathcal{O}(N^2)$};
  \end{scope}
  % --- destra: Q (K^T V) -----------------------------------------------------
  \begin{scope}[xshift=92mm]
    \draw[mat, fill=SapRossoChiaro] (0,20) rectangle (24,24);     \node at (12,17) {$\hat K\T$};
    \node at (28,12) {$\cdot$};
    \draw[mat, fill=SapGrigioChiaro, draw=SapGrigio] (32,0) rectangle (38,24); \node at (35,-3) {$V$};
    \node at (42,12) {$=$};
    \draw[mat, fill=SapRossoTenue] (46,18) rectangle (52,24);   \node[align=center] at (49,13) {$d_k\times C$};
    \node at (57,12) {$\Rightarrow$};
    \draw[mat, fill=SapRossoChiaro] (62,0) rectangle (66,24);   \node at (64,-3) {$\hat Q$};
    \node at (69,12) {$\cdot$};
    \draw[mat, fill=SapRossoTenue] (72,18) rectangle (78,24);
    \node[text=SapRossoScuro, font=\footnotesize\bfseries] at (39,31) {$\hat Q\,(\hat K\T V)$: memoria e calcolo $\mathcal{O}(N)$};
  \end{scope}
\end{tikzpicture}
}
  \caption[Ordine dei prodotti nell'attenzione]{L'ordine dei prodotti determina il costo. A
    sinistra la forma dot-product materializza la matrice $N\times N$. A destra la forma lineare
    calcola prima $\hat K\T V\in\R^{d_k\times C}$, comune a tutte le query.}
  \label{fig:associativita}
\end{figure}

\subsection{Il Linear Attention Mechanism}
\label{sec:attenzione-lineare}

Il LAM di Li et al.~\cite{li2020lam} ottiene una fattorizzazione del tipo
\eqref{eq:attenzione-lineare-generale} sviluppando l'esponenziale al primo ordine,
$e^{\vect{q}\T\vect{k}}\approx1+\vect{q}\T\vect{k}$. Lo sviluppo da solo non garantisce pesi non
negativi, perché $\vect{q}\T\vect{k}$ può essere minore di $-1$. Per questo query e key vengono
normalizzate in norma $\ell_2$, $\hatq_i=\vect{q}_i/\norm{\vect{q}_i}_2$ e
$\hatk_j=\vect{k}_j/\norm{\vect{k}_j}_2$, e si usa
\begin{equation}
  \mathrm{sim}(\vect{q}_i,\vect{k}_j)=1+\hatq_i\T\hatk_j=1+\cos\vartheta_{ij}\in[0,2].
\end{equation}
Sostituendo nella~\eqref{eq:attenzione-nucleo} si ottiene
\begin{equation}
  D_i=\frac{\sum_{j}\vect{v}_j+\Bigl(\sum_{j}\hatk_j\vect{v}_j\T\Bigr)\T\hatq_i}
           {N+\hatq_i\T\sum_{j}\hatk_j},
  \label{eq:lam}
\end{equation}
che corrisponde alla~\eqref{eq:attenzione-lineare-generale} con
$\phi(\vect{q})=(1,\hat{\vect{q}})$ e $\psi(\vect{k})=(1,\hat{\vect{k}})$, quindi con $m=d_k+1$.
In forma matriciale, con $\hat Q$ e $\hat K$ normalizzate riga per riga,
$D=\bigl(\vect{1}\vect{1}\T V+\hat Q(\hat K\T V)\bigr)\oslash\bigl(N\vect{1}+\hat Q\hat K\T\vect{1}\bigr)$,
dove $\oslash$ divide ogni riga per il corrispondente elemento del vettore. Il calcolo richiede
$\bigO{N d_k C}$ operazioni e memoria $\bigO{N(d_k+C)+d_kC}$. Nell'implementazione si aggiunge
una piccola costante $\epsilon$ alla somma delle key per stabilità numerica.

\begin{osservazione}[Contrasto limitato del LAM]
\label{oss:contrasto}
Sia $\bar{\vect{m}}=\frac1N\sum_j\hatk_j$ la media delle key normalizzate, con
$\norm{\bar{\vect{m}}}_2\le1$. Il peso che la posizione $i$ assegna alla posizione $j$ è
\[
  \alpha_{ij}=\frac{1+\hatq_i\T\hatk_j}{N\,(1+\hatq_i\T\bar{\vect{m}})}
  \;\leq\;\frac{2}{N\,(1-\norm{\bar{\vect{m}}}_2)} .
\]
Se le key non sono tutte allineate ($\norm{\bar{\vect{m}}}_2$ lontana da 1), nessuna posizione
riceve più di una piccola frazione costante del peso medio $1/N$. La softmax, al contrario, può
concentrare quasi tutto il peso su una sola posizione. Il LAM si comporta quindi come
un'\emph{aggregazione di contesto globale modulata} dalla somiglianza coseno più che come un
meccanismo di selezione. Il caso $\norm{\bar{\vect{m}}}_2\to1$ con
$\hatq_i\approx-\bar{\vect{m}}$ è quello in cui il denominatore si annulla e serve
$\epsilon$.
\end{osservazione}

L'osservazione suggerisce che il contributo del LAM vada cercato nella coerenza su larga scala
(per esempio nell'uniformità delle grandi superfici) più che nella precisione dei contorni.
\Segnaposto{Verificare l'osservazione sui modelli addestrati: distribuzione di
$\norm{\bar{\vect{m}}}_2$ e di $\max_j\alpha_{ij}$ per stadio, da confrontare con le classi in
cui \RQ{1} mostra differenze.}

L'uscita del blocco è residuale, $\vect{X}'=\vect{X}+\gamma\,D(\vect{X})$, con un guadagno
scalare $\gamma$ appreso e \textbf{inizializzato a zero}. All'inizio dell'addestramento il
blocco è quindi l'identità e il modello parte dal comportamento della rete senza attenzione.

\subsection{Attenzione sui canali}
\label{sec:cam}

Sui canali il costo quadratico non è un problema, perché $C\ll N$. Seguendo
DANet~\cite{fu2019danet}, il modulo CAM calcola la matrice di affinità tra canali
$E=\vect{X}\T\vect{X}\in\R^{C\times C}$ e i pesi
$A=\softmax_{\mathrm{riga}}\bigl(\max_{\mathrm{riga}}(E)-E\bigr)$, poi produce
$\vect{X}'=\vect{X}+\beta\,\vect{X}A\T$ con $\beta$ inizializzato a zero. Il costo è
$\bigO{NC^2}$.

\begin{osservazione}
La softmax è invariante per l'aggiunta di una costante a ciascuna riga, quindi
$\softmax_{\mathrm{riga}}(\max_{\mathrm{riga}}(E)-E)=\softmax_{\mathrm{riga}}(-E)$. La
sottrazione del massimo è solo un accorgimento di stabilità numerica, e l'effetto netto è di
dare più peso ai canali \emph{meno} correlati con il canale considerato.
\end{osservazione}

\subsection{L'architettura MAResU-Net}
\label{sec:maresunet}

\maresunet{}~\cite{li2022maresunet} instanzia la~\eqref{eq:decoder} con un blocco di attenzione
su ciascuna delle quattro scale (\cref{fig:maresunet,fig:pamcam}). Il blocco applica una
convoluzione $3\times3$, poi in parallelo il PAM lineare~\eqref{eq:lam} e il CAM, ciascuno
seguito da dropout spaziale e convoluzione $1\times1$. Le due uscite vengono sommate e
proiettate con una convoluzione $1\times1$. Il decoder è
\begin{align}
  \vect{S}_l &= \mathrm{PAMCAM}_l(\vect{E}_l),\quad l=1,\dots,4, &
  \vect{B} &= \mathrm{Neck}(\vect{S}_4), \nonumber\\
  \vect{D}_3 &= \mathrm{Dec}_4(\vect{B})+\vect{S}_3, &
  \vect{D}_2 &= \mathrm{Dec}_3(\vect{D}_3)+\vect{S}_2, \label{eq:maresunet}\\
  \vect{D}_1 &= \mathrm{Dec}_2(\vect{D}_2)+\vect{S}_1, &
  \vect{z} &= \mathrm{Head}\bigl(\mathrm{Dec}_1(\vect{D}_1)\bigr). \nonumber
\end{align}
Ogni $\mathrm{Dec}_l$ è la sequenza convoluzione $1\times1$ (riduzione dei canali a un quarto),
convoluzione trasposta $3\times3$ con stride 2 e convoluzione $1\times1$, ciascuna con batch
normalization~\cite{ioffe2015batchnorm} e ReLU. La testa riporta la risoluzione piena con una
convoluzione trasposta $4\times4$ con stride 2 seguita da due convoluzioni $3\times3$. Il
\emph{neck} è l'identità nella \maresunet{} e un quantizzatore vettoriale nella \maresunetvq{}
(\cref{sec:vq}).

\begin{figure}[p]
  \centering
  \resizebox{\linewidth}{!}{% Architettura MAResU-Net: encoder ResNet, blocchi PAM-CAM sulle connessioni di salto,
% neck opzionale (identità o VQ), decoder con somme e testa a risoluzione piena.
\begin{tikzpicture}[
    e/.style={enc, minimum width=2.3cm},
    a/.style={attn, minimum width=2.3cm},
    d/.style={dec, minimum width=2.3cm},
    somma/.style={circle, draw=SapRossoScuro, fill=white, inner sep=0pt, minimum size=4.2mm,
                  font=\footnotesize, line width=0.5pt},
  ]
  % --- livelli (stride) ---------------------------------------------------
  \def\yA{-2.0}\def\yB{-3.6}\def\yC{-5.2}\def\yD{-6.8}

  % --- encoder --------------------------------------------------------------
  \node[dato, minimum width=2.3cm] (in) at (0,0.1) {Immagine\\{\scriptsize $H\times W\times3$ (o 4 con nDSM)}};
  \node[e] (stem) at (0,-0.95) {conv $7{\times}7$, s2 + max-pool};
  \node[e] (e1) at (0,\yA) {\texttt{layer1}\\{\scriptsize 64 canali}};
  \node[e] (e2) at (0,\yB) {\texttt{layer2}\\{\scriptsize 128 canali}};
  \node[e] (e3) at (0,\yC) {\texttt{layer3}\\{\scriptsize 256 canali}};
  \node[e] (e4) at (0,\yD) {\texttt{layer4}\\{\scriptsize 512 canali}};
  \draw[flusso] (in) -- (stem);
  \draw[flusso] (stem) -- (e1);
  \draw[flusso] (e1) -- (e2);
  \draw[flusso] (e2) -- (e3);
  \draw[flusso] (e3) -- (e4);

  % --- etichette di stride ----------------------------------------------------
  \foreach \y/\s in {\yA/{1/4}, \yB/{1/8}, \yC/{1/16}, \yD/{1/32}}{
    \node[nota, anchor=east] at (-1.45,\y) {$\s$};
  }
  \node[nota, anchor=east, font=\scriptsize\bfseries] at (-1.45,-1.45) {scala};

  % --- attenzione sulle connessioni di salto ---------------------------------
  \node[a] (a1) at (3.9,\yA) {PAM lineare $\parallel$ CAM};
  \node[a] (a2) at (3.9,\yB) {PAM lineare $\parallel$ CAM};
  \node[a] (a3) at (3.9,\yC) {PAM lineare $\parallel$ CAM};
  \node[a] (a4) at (3.9,\yD) {PAM lineare $\parallel$ CAM};
  \foreach \i in {1,...,4}{ \draw[flusso] (e\i) -- (a\i); }

  % --- neck ------------------------------------------------------------------
  \node[quant, minimum width=1.9cm] (neck) at (6.9,\yD) {Neck\\{\scriptsize identità o VQ}};
  \draw[flusso] (a4) -- (neck);

  % --- decoder ----------------------------------------------------------------
  \node[d] (d4) at (9.6,{(\yC+\yD)/2}) {Dec$_4$ \scriptsize($\uparrow2$)};
  \node[somma] (s3) at (9.6,\yC) {$+$};
  \node[d] (d3) at (9.6,{(\yB+\yC)/2}) {Dec$_3$ \scriptsize($\uparrow2$)};
  \node[somma] (s2) at (9.6,\yB) {$+$};
  \node[d] (d2) at (9.6,{(\yA+\yB)/2}) {Dec$_2$ \scriptsize($\uparrow2$)};
  \node[somma] (s1) at (9.6,\yA) {$+$};
  \node[d] (d1) at (9.6,-0.95) {Dec$_1$ \scriptsize($\uparrow2$, 1/2)};
  \node[d, fill=SapRossoVelo, text width=2.5cm] (head) at (9.6,0.25) {Testa\\{\scriptsize ConvT $4{\times}4$ s2,\\2 conv $3{\times}3$}};
  \node[dato, anchor=west, minimum width=1.6cm] (out) at (11.8,0.25) {Logit\\{\scriptsize $H\times W\times6$}};

  \draw[flusso] (neck) -| (d4);
  \draw[flusso] (d4) -- (s3);
  \draw[flusso] (s3) -- (d3);
  \draw[flusso] (d3) -- (s2);
  \draw[flusso] (s2) -- (d2);
  \draw[flusso] (d2) -- (s1);
  \draw[flusso] (s1) -- (d1);
  \draw[flusso] (d1) -- (head);
  \draw[flusso] (head) -- (out);

  % connessioni di salto (dopo l'attenzione)
  \draw[skipc] (a3) -- (s3);
  \draw[skipc] (a2) -- (s2);
  \draw[skipc] (a1) -- (s1);

  % --- raggruppamenti -----------------------------------------------------------
  \begin{scope}[on background layer]
    \node[gruppo, fit=(stem)(e4), label={[nota]below:encoder (ResNet-18/34)}] {};
    \node[gruppo, fit=(d1)(d4)(head), label={[nota]above:decoder}] {};
  \end{scope}
\end{tikzpicture}
}
  \caption[Architettura MAResU-Net]{Architettura \maresunet{}. Le mappe dell'encoder a stride
    4--32 passano da un blocco di attenzione (PAM lineare e CAM in parallelo) prima di essere
    sommate alle uscite del decoder. Il neck sul collo di bottiglia è l'identità oppure un
    quantizzatore vettoriale. Le frecce rosse sono le connessioni di salto.}
  \label{fig:maresunet}
  \vspace{1.2em}
  \resizebox{\linewidth}{!}{% Blocco di attenzione di MAResU-Net: conv 3x3, poi PAM lineare e CAM in parallelo,
% somma e proiezione 1x1. Entrambi i rami sono residuali con guadagno inizializzato a 0.
\begin{tikzpicture}[node distance=3.5mm and 6mm,
    b/.style={blocco, fill=white, draw=SapGrigio},
    somma/.style={circle, draw=SapRossoScuro, fill=white, inner sep=0pt, minimum size=4.2mm,
                  font=\footnotesize},
  ]
  \node[dato] (x) {$\vect{X}_l$\\{\scriptsize $h\times w\times C$}};
  \node[b, right=of x] (c3) {conv $3{\times}3$\\ReLU};
  \node[attn, above right=2mm and 7mm of c3, minimum width=2.6cm] (pam)
       {PAM lineare\\{\scriptsize $\vect{X}+\gamma\,\mathrm{LAM}(\vect{X})$}};
  \node[attn, below right=2mm and 7mm of c3, minimum width=2.6cm] (cam)
       {CAM\\{\scriptsize $\vect{X}+\beta\,\mathrm{CA}(\vect{X})$}};
  \node[b, right=of pam] (p1) {drop\\conv $1{\times}1$};
  \node[b, right=of cam] (c1) {drop\\conv $1{\times}1$};
  \node[somma, right=12mm of $(p1)!0.5!(c1)$] (s) {$+$};
  \node[b, right=of s] (o1) {drop\\conv $1{\times}1$};
  \node[dato, right=of o1] (y) {$\vect{S}_l$};
  \draw[flusso] (x) -- (c3);
  \draw[flusso] (c3.east) -- ++(3mm,0) |- (pam.west);
  \draw[flusso] (c3.east) -- ++(3mm,0) |- (cam.west);
  \draw[flusso] (pam) -- (p1);
  \draw[flusso] (cam) -- (c1);
  \draw[flusso] (p1.east) -| (s.north);
  \draw[flusso] (c1.east) -| (s.south);
  \draw[flusso] (s) -- (o1);
  \draw[flusso] (o1) -- (y);
  \node[nota, below=1mm of cam] {$\gamma=\beta=0$ all'inizializzazione};
\end{tikzpicture}
}
  \caption[Blocco di attenzione PAM-CAM]{Blocco di attenzione di \maresunet{}. Entrambi i rami
    sono residuali con guadagno iniziale nullo: all'inizio dell'addestramento il blocco calcola
    soltanto le convoluzioni.}
  \label{fig:pamcam}
\end{figure}

\begin{table}[t]
  \centering
  \small
  \caption[Costo dei moduli di attenzione]{Costo asintotico dei moduli di attenzione su una
    mappa con $N$ posizioni e $C$ canali ($d_k=C/8$ nel PAM).}
  \label{tab:costo-attenzione}
  \begin{tabular}{|l|c|c|l|}
    \hline
    \textbf{Modulo} & \textbf{Calcolo} & \textbf{Memoria aggiuntiva} & \textbf{Uso} \\
    \hline
    PAM dot-product & $\bigO{N^2(d_k+C)}$ & $\bigO{N^2}$ & solo mappe piccole \\
    PAM lineare (LAM) & $\bigO{N d_k C}$ & $\bigO{N(d_k+C)+d_kC}$ & tutti gli stadi \\
    CAM & $\bigO{NC^2}$ & $\bigO{C^2}$ & tutti gli stadi \\
    \hline
  \end{tabular}
\end{table}

% ---------------------------------------------------------------------
\section{Quantizzazione vettoriale}
\label{sec:vq}

\subsection{Codebook e assegnazione}

Un quantizzatore vettoriale~\cite{vandenoord2017vqvae} mantiene un \emph{codebook}
$\mathcal{E}=\{\vect{e}_1,\dots,\vect{e}_K\}\subset\R^{D}$ e sostituisce ogni vettore
$\vect{z}_e(u)\in\R^{D}$ prodotto dall'encoder con il codice più vicino:
\begin{equation}
  k^\star(u)=\argmin_{k}\norm{\vect{z}_e(u)-\vect{e}_k}_2^2,
  \qquad
  \vect{z}_q(u)=\vect{e}_{k^\star(u)} .
  \label{eq:vq-assegnazione}
\end{equation}
Lo spazio delle caratteristiche viene così partizionato nelle celle di Voronoi dei codici
(\cref{fig:vq}b), e il decoder riceve uno dei soli $K$ valori possibili per ogni posizione. La
distanza si calcola senza materializzare differenze, come
$\norm{\vect{z}}^2-2\vect{z}\T\vect{e}_k+\norm{\vect{e}_k}^2$.

\begin{figure}[t]
  \centering
  \resizebox{\linewidth}{!}{% Quantizzazione vettoriale: (a) flusso in avanti e gradiente straight-through;
% (b) geometria: assegnazione al codice più vicino, commitment e aggiornamento EMA.
\begin{tikzpicture}[node distance=4mm and 7mm]
  % ---------------- (a) flusso -------------------------------------------------
  \begin{scope}[local bounding box=A]
    \node[enc] (enc) {Encoder\\$f_\theta$};
    \node[dato, right=of enc] (ze) {$\vect{z}_e(u)\in\R^{D}$};
    \node[quant, right=of ze] (nn) {$k^\star=\argmin_k \norm{\vect{z}_e-\vect{e}_k}_2$};
    \node[dato, right=of nn] (zq) {$\vect{z}_q(u)=\vect{e}_{k^\star}$};
    \node[dec, right=of zq] (dec) {Decoder\\$g_\psi$};
    \node[blocco, fill=SapRossoVelo, below=6mm of nn, minimum width=3.6cm] (cb)
         {codebook $\{\vect{e}_1,\dots,\vect{e}_K\}\subset\R^{D}$};
    \draw[flusso] (enc) -- (ze);
    \draw[flusso] (ze) -- (nn);
    \draw[flusso] (nn) -- (zq);
    \draw[flusso] (zq) -- (dec);
    \draw[flusso] (cb) -- (nn);
    % gradiente straight-through
    \draw[grad] (dec.north) to[bend right=22]
      node[midway, above, font=\scriptsize, text=SapRosso]
      {$\nabla_{\vect{z}_q}\mathcal{L}$ copiato su $\vect{z}_e$ (straight-through)} (ze.north);
    \node[nota, anchor=north west] at ($(enc.south west)+(0,-3mm)$)
      {\textcolor{ColFreccia}{$\longrightarrow$} in avanti \quad
       \textcolor{SapRosso}{$\dashrightarrow$} gradiente};
  \end{scope}
  \node[font=\small\bfseries, anchor=south west] at ($(A.north west)+(0,1mm)$) {(a)};

  % ---------------- (b) geometria ----------------------------------------------
  \begin{scope}[shift={($(A.south west)+(1.0cm,-5.2cm)$)}, x=1.15cm, y=1.15cm, local bounding box=B]
    % cella del codice e_2: vettori assegnati e loro centroide
    \fill[SapRossoVelo] (2,1.6) circle (1.15);
    \foreach \q in {(1.35,1.15),(2.55,1.05),(1.6,2.35),(2.65,2.15),(1.2,1.85)}{
      \draw[SapGrigio, line width=0.5pt] \q ++(-2pt,-2pt) -- ++(4pt,4pt) \q ++(-2pt,2pt) -- ++(4pt,-4pt);
    }
    \coordinate (c) at (1.87,1.71);
    \draw[SapGrigio, very thick] (c) circle (2.4pt);
    % codici
    \foreach \p/\n/\a in {{(0,0)}/1/below, {(2,1.6)}/2/below right, {(4.6,0.2)}/3/below,
                          {(4.0,3.3)}/4/above, {(-0.3,2.9)}/5/above, {(6.6,3.1)}/6/above}{
      \fill[SapRossoScuro] \p circle (1.8pt) node[\a, font=\scriptsize] {$\vect{e}_{\n}$};
    }
    % vettore dell'encoder e commitment
    \coordinate (z) at (3.55,0.55);
    \draw[black, line width=0.7pt] (z) ++(-2.5pt,-2.5pt) -- ++(5pt,5pt) (z) ++(-2.5pt,2.5pt) -- ++(5pt,-5pt);
    \node[font=\scriptsize, anchor=north west] at (3.6,0.5) {$\vect{z}_e$};
    \draw[->, SapRosso, thick] (z) -- ($(z)!0.5!(2,1.6)$);
    \node[font=\scriptsize, text=SapRosso, anchor=south west] at (3.05,1.0) {commitment ($\beta$)};
    % aggiornamento EMA del codice
    \draw[->, SapRossoScuro, thick] (2,1.6) -- ($(2,1.6)!0.85!(c)$);
    % didascalie interne
    \node[font=\scriptsize, text=SapGrigio, align=center, anchor=south] at (2,2.8)
      {vettori assegnati a $\vect{e}_2$ ($\times$)\\e loro centroide ($\circ$)};
    \node[font=\scriptsize, text=SapRossoScuro, align=left, anchor=west] at (4.75,2.15)
      {EMA: $\vect{e}_2$ si sposta verso\\il centroide della sua cella};
    \draw[SapRossoScuro, densely dotted] (4.72,2.15) -- (2.12,1.66);
  \end{scope}
  \node[font=\small\bfseries, anchor=south west] at ($(B.north west)+(-1.0cm,1mm)$) {(b)};
\end{tikzpicture}
}
  \caption[Quantizzazione vettoriale]{Quantizzazione vettoriale. (a)~Flusso in avanti e
    gradiente \emph{straight-through}: il gradiente rispetto a $\vect{z}_q$ viene copiato su
    $\vect{z}_e$. (b)~Geometria nello spazio delle caratteristiche: il vettore $\vect{z}_e$ è
    assegnato al codice più vicino, il termine di commitment lo avvicina al codice e
    l'aggiornamento EMA sposta il codice verso il centroide dei vettori assegnati.}
  \label{fig:vq}
\end{figure}

\subsection{Gradiente straight-through e perdita}

L'$\argmin$ ha derivata nulla quasi ovunque. Lo stimatore
\emph{straight-through}~\cite{bengio2013straight} usa nel passo in avanti il valore quantizzato
e nel passo all'indietro l'identità:
\begin{equation}
  \tilde{\vect{z}}=\vect{z}_e+\sg\!\left[\vect{z}_q-\vect{z}_e\right],
  \qquad
  \tilde{\vect{z}}=\vect{z}_q,\quad
  \frac{\partial\tilde{\vect{z}}}{\partial\vect{z}_e}=I,
  \label{eq:straight-through}
\end{equation}
dove $\sg[\cdot]$ è l'operatore \emph{stop-gradient} (identità in avanti, derivata nulla). Il
codebook non riceve gradiente dal decoder e va appreso con un termine apposito. Si aggiunge
inoltre un termine di \emph{commitment} che impedisce alle uscite dell'encoder di allontanarsi
indefinitamente dai codici:
\begin{equation}
  \mathcal{L}_{\mathrm{VQ}}=
  \underbrace{\norm{\sg[\vect{z}_e]-\vect{z}_q}_2^2}_{\text{codebook}}
  +\beta\,\underbrace{\norm{\vect{z}_e-\sg[\vect{z}_q]}_2^2}_{\text{commitment}},
  \qquad
  \mathcal{L}=\mathcal{L}_{\mathrm{CE}}+\lambda\,\mathcal{L}_{\mathrm{VQ}},
  \label{eq:vq-loss}
\end{equation}
con le norme mediate sulle posizioni e sulle componenti, $\beta=0{,}25$ e $\lambda$ peso della
perdita ausiliaria.

\subsection{Aggiornamento EMA come K-means online}

In alternativa al termine \enquote{codebook} della~\eqref{eq:vq-loss}, i codici possono essere
aggiornati con medie mobili esponenziali~\cite{vandenoord2017vqvae,razavi2019vqvae2}. Sia, per
un batch, $n_k=\sum_u\Ind[k^\star(u)=k]$ il numero di vettori assegnati al codice $k$ e
$\vect{s}_k=\sum_u\Ind[k^\star(u)=k]\,\vect{z}_e(u)$ la loro somma. Con decadimento
$\gamma\in(0,1)$:
\begin{equation}
  N_k\leftarrow\gamma N_k+(1-\gamma)\,n_k,\qquad
  \vect{M}_k\leftarrow\gamma\vect{M}_k+(1-\gamma)\,\vect{s}_k,\qquad
  \vect{e}_k\leftarrow\vect{M}_k/\tilde N_k ,
  \label{eq:vq-ema}
\end{equation}
dove $\tilde N_k=\frac{N_k+\varepsilon}{\sum_j N_j+K\varepsilon}\sum_jN_j$ è una versione
regolarizzata (smoothing di Laplace) che evita divisioni per zero. Con l'EMA la perdita si
riduce al solo commitment, $\mathcal{L}_{\mathrm{VQ}}=\beta\norm{\vect{z}_e-\sg[\vect{z}_q]}^2$.

\begin{proposizione}[EMA e K-means]
\label{prop:ema-kmeans}
Trascurando lo smoothing e l'inizializzazione, dopo $t$ passi l'aggiornamento~\eqref{eq:vq-ema}
restituisce
\[
  \vect{e}_k^{(t)}=\frac{\sum_{\tau\le t}\gamma^{t-\tau}\vect{s}_k^{(\tau)}}
                        {\sum_{\tau\le t}\gamma^{t-\tau}n_k^{(\tau)}}
  =\argmin_{\vect{e}}\;\sum_{\tau\le t}\gamma^{t-\tau}\sum_{u}\Ind\!\left[k^{\star(\tau)}(u)=k\right]
   \norm{\vect{z}_e^{(\tau)}(u)-\vect{e}}_2^2 ,
\]
cioè il centroide dei vettori assegnati al codice $k$, pesati con uno sconto esponenziale
sull'età del batch. È un passo di Lloyd~\cite{lloyd1982kmeans} eseguito online.
\end{proposizione}
\begin{proof}
Srotolando le ricorrenze, $\vect{M}_k^{(t)}=(1-\gamma)\sum_{\tau}\gamma^{t-\tau}\vect{s}_k^{(\tau)}$ e
$N_k^{(t)}=(1-\gamma)\sum_\tau\gamma^{t-\tau}n_k^{(\tau)}$. Il fattore $(1-\gamma)$ si semplifica
nel rapporto. La funzione a destra è quadratica e convessa in $\vect{e}$, e annullandone il
gradiente,
$-2\sum_\tau\gamma^{t-\tau}\bigl(\vect{s}_k^{(\tau)}-n_k^{(\tau)}\vect{e}\bigr)=\vect{0}$, si
ottiene lo stesso rapporto.
\end{proof}

La \cref{prop:ema-kmeans} spiega perché lo stato EMA va trattato come stato persistente e non
come parametro ricreato a ogni passo, come accadeva nel 2022. Spiega anche il collasso tipico del
codebook: un codice che non riceve assegnazioni non si muove più, perché ha $n_k=0$ in ogni
batch, e diventa un \emph{codice morto}. Le contromisure adottate sono due
(\cref{alg:vq-ema}): l'\textbf{inizializzazione dai dati}, che pone i codici su vettori estratti
dal primo batch invece che in un intorno dell'origine lontano dalle attivazioni post-ReLU, e il
\textbf{riavvio dei codici morti}~\cite{lancucki2020robust}, che riposiziona su un vettore del
batch ogni codice con $N_k$ sotto una soglia.

\begin{algorithm}[t]
  \caption{Passo di addestramento del quantizzatore con EMA e riavvio dei codici morti}
  \label{alg:vq-ema}
  \begin{algorithmic}[1]
    \Require uscite dell'encoder $Z=\{\vect{z}_e(u)\}_{u}$, codebook $\{\vect{e}_k\}$, stato $\{N_k,\vect{M}_k\}$, decadimento $\gamma$, soglia $\tau_{\mathrm{morto}}$
    \If{primo passo}
      \State $\vect{e}_k\gets$ campione casuale di $Z$ (con piccolo rumore), $\vect{M}_k\gets\vect{e}_k$, $N_k\gets1$ per ogni $k$
    \EndIf
    \State $k^\star(u)\gets\argmin_k\norm{\vect{z}_e(u)-\vect{e}_k}^2$; \quad $\vect{z}_q(u)\gets\vect{e}_{k^\star(u)}$
    \State $\mathcal{L}_{\mathrm{VQ}}\gets\beta\cdot\mathrm{media}_u\norm{\vect{z}_e(u)-\sg[\vect{z}_q(u)]}^2$
    \State $N_k\gets\gamma N_k+(1-\gamma)n_k$; \quad $\vect{M}_k\gets\gamma\vect{M}_k+(1-\gamma)\vect{s}_k$
    \ForEach{$k$ con $N_k<\tau_{\mathrm{morto}}$}
      \State $\vect{M}_k\gets$ campione casuale di $Z$; \quad $N_k\gets1$ \Comment{riavvio}
    \EndFor
    \State $\vect{e}_k\gets\vect{M}_k/\tilde N_k$
    \State \Return $\vect{z}_e+\sg[\vect{z}_q-\vect{z}_e]$, $\mathcal{L}_{\mathrm{VQ}}$
  \end{algorithmic}
\end{algorithm}

\subsection{Diagnostica: perplexity e uso del codebook}

Sia $\bar p_k=n_k/\sum_jn_j$ la frequenza di assegnazione del codice $k$ in un batch. Si
definiscono
\begin{equation}
  \PPL=\exp\Bigl(-\sum_{k}\bar p_k\log\bar p_k\Bigr)\in[1,K],
  \qquad
  \mathrm{uso}=\frac{\abs{\{k:\bar p_k>0\}}}{K}\in\left[\tfrac1K,1\right].
  \label{eq:perplexity}
\end{equation}
La perplexity è il numero \emph{effettivo} di codici usati: vale 1 se tutti i vettori finiscono
sullo stesso codice e $K$ se l'uso è uniforme. Un codebook \enquote{sano}, nel senso della
\RQ{2}, ha perplexity alta e uso vicino a 1 durante tutto l'addestramento.

\subsection{Capacità del collo di bottiglia e aggiramento}
\label{sec:vq-aggiramento}

La quantizzazione limita l'informazione che attraversa il collo di bottiglia. Con $h\times w$
posizioni e $K$ codici,
\begin{equation}
  I(\vect{x};\vect{z}_q)\;\le\;\Hent(\vect{z}_q)\;\le\;hw\log_2K\ \text{bit}.
  \label{eq:capacita}
\end{equation}
Per un ritaglio di $512\times512$ pixel il collo di bottiglia di \maresunet{} è a stride 32
($16\times16$ posizioni). Con $K=512$ porta al più $256\cdot9=\num{2304}$ bit, cioè 288 byte per
ritaglio. La sola connessione di salto a stride 4 trasporta invece
$128\cdot128\cdot64\approx10^6$ valori reali. Il collo di bottiglia vincola quindi solo il
percorso più stretto della rete, che non è l'unico. Ne segue il risultato seguente.

\begin{proposizione}[Aggiramento del collo di bottiglia]
\label{prop:aggiramento}
Si consideri la~\eqref{eq:maresunet} con il quantizzatore solo sul neck e si supponga che la
perdita di segmentazione $\mathcal{L}_{\mathrm{CE}}$ ammetta un minimo globale con un decoder
che non dipende da $\vect{B}$. Allora la perdita totale~\eqref{eq:vq-loss} ammette un minimo
globale in cui il codebook usa un solo codice ($\PPL=1$).
\end{proposizione}
\begin{proof}
Si parta dal minimo di $\mathcal{L}_{\mathrm{CE}}$ con decoder indipendente da $\vect{B}$. Le
mappe $\vect{S}_1,\vect{S}_2,\vect{S}_3$ non dipendono dallo stadio $\vect{E}_4$. Si possono
quindi modificare i parametri dell'ultimo stadio (per esempio azzerando il guadagno dell'ultima
normalizzazione) in modo che $\vect{E}_4$, e quindi $\vect{z}_e$, sia costante e pari a
$\vect{c}$, senza cambiare $\mathcal{L}_{\mathrm{CE}}$. Ponendo $\vect{e}_1=\vect{c}$, il
commitment è nullo e tutti i vettori sono assegnati a $\vect{e}_1$. Entrambi i termini, che sono
non negativi, sono al loro minimo.
\end{proof}

La proposizione non afferma che l'ottimizzazione raggiunga quella soluzione. Afferma che
\emph{nulla nell'obiettivo la penalizza}: con le sole connessioni di salto libere, un risultato
negativo della \maresunetvq{} non dice nulla sull'utilità della quantizzazione. Per questo la
\RQ{2} misura la perplexity durante l'addestramento e prevede una variante in cui anche le
connessioni di salto sono quantizzate, così che ogni percorso dall'encoder al decoder passi da
un codebook e l'ipotesi della \cref{prop:aggiramento} cada.

% ---------------------------------------------------------------------
\section{Auto-distillazione senza etichette}
\label{sec:dino}

\subsection{Student, teacher e distribuzioni}

DINO~\cite{caron2021dino} addestra uno \emph{student} $g_s=h_s\circ\GAP\circ f_s$ a riprodurre
l'uscita di un \emph{teacher} $g_t=h_t\circ\GAP\circ f_t$ con la stessa architettura. Qui $f$ è
l'encoder (lo stesso che verrà usato in \maresunet{}), $\GAP$ la media globale sulle posizioni e
$h$ una testa di proiezione: un MLP a tre strati (2048, 2048, 256 unità con GELU),
normalizzazione $\ell_2$ e uno strato lineare a peso normalizzato verso $K=\num{16384}$
\enquote{prototipi}. Le uscite diventano distribuzioni con due temperature:
\begin{equation}
  P_s(\vect{x})=\softmax\!\left(\frac{g_s(\vect{x})}{\tau_s}\right),\qquad
  P_t(\vect{x})=\softmax\!\left(\frac{g_t(\vect{x})-\vect{c}}{\tau_t}\right),
  \label{eq:dino-distribuzioni}
\end{equation}
dove $\vect{c}\in\R^K$ è il \emph{centro} e $\tau_t<\tau_s$ rende la distribuzione del teacher
più concentrata (\emph{sharpening}).

\subsection{Multi-crop e perdita}

Da ogni immagine si estraggono due viste globali $\vect{x}^{g_1},\vect{x}^{g_2}$ ($224\times224$,
scala tra il 14\% e il 100\% dell'area) e $M$ viste locali $\vect{x}^{l_1},\dots,\vect{x}^{l_M}$
($96\times96$, scala tra il 5\% e il 40\%), con trasformazioni cromatiche e geometriche
casuali. Il teacher vede solo le viste globali, lo student tutte. Per un'immagine la perdita è
\begin{equation}
  \ell(\vect{x})=\frac{1}{\abs{\mathcal{P}}}\sum_{(a,b)\in\mathcal{P}}
  \Hent\bigl(P_t(\vect{x}^{a}),P_s(\vect{x}^{b})\bigr),
  \qquad
  \mathcal{P}=\bigl\{(a,b):a\in\{g_1,g_2\},\ b\in V,\ b\neq a\bigr\},
  \label{eq:dino-loss}
\end{equation}
con $V$ l'insieme di tutte le viste, $\Hent(p,q)=-\sum_kp_k\log q_k$ e
$\abs{\mathcal{P}}=2(M+1)$. Sul batch si prende la media
$\mathcal{L}=\frac1B\sum_i\ell(\vect{x}_i)$. La perdita incoraggia corrispondenze
\enquote{da locale a globale}: lo student deve dedurre da una porzione dell'immagine la
distribuzione che il teacher assegna a una vista ampia.

\begin{osservazione}
La normalizzazione~\eqref{eq:dino-distribuzioni} è \emph{per campione}: ogni immagine ha la sua
distribuzione su $K$ prototipi. Calcolare la softmax sull'intero batch appiattito, come avveniva
nel codice del 2022, produce un'unica distribuzione su $B\cdot K$ elementi (lì addirittura sui
pixel dell'uscita densa del decoder). Le immagini del batch competono allora per la stessa massa
di probabilità e l'obiettivo non è più quello di DINO.
\end{osservazione}

\begin{figure}[t]
  \centering
  \resizebox{\linewidth}{!}{% Auto-distillazione DINO con multi-crop: teacher sulle viste globali, student su tutte;
% centering e sharpening sul teacher, stop-gradient, teacher come EMA dello student.
\begin{tikzpicture}[node distance=4mm and 6mm,
    vista/.style={dato, minimum width=1.2cm, minimum height=6mm, font=\scriptsize},
  ]
  \node[dato, minimum height=12mm] (x) {immagine\\$\vect{x}$};

  \node[vista, above right=3mm and 9mm of x] (g) {$\vect{x}^{g}_{1}, \vect{x}^{g}_{2}$\\$224^2$};
  \node[vista, below right=3mm and 9mm of x] (l) {$\vect{x}^{l}_{1..M}$\\$96^2$};
  \draw[flusso] (x.east) -- ++(3mm,0) |- (g.west);
  \draw[flusso] (x.east) -- ++(3mm,0) |- (l.west);
  \node[nota, above=0mm of g] {viste globali};
  \node[nota, below=0mm of l] {viste locali};

  % teacher
  \node[enc, right=14mm of g, minimum width=2.5cm, fill=SapGrigioChiaro, draw=SapGrigio] (t)
       {teacher $h_t\circ f_t$\\{\scriptsize $\theta_t$, nessun gradiente}};
  \node[blocco, fill=white, right=of t] (tc) {$\softmax\!\big((\cdot-\vect{c})/\tau_t\big)$};
  \node[dato, right=of tc] (pt) {$P_t$};

  % student
  \node[enc, right=14mm of l, minimum width=2.5cm] (s) {student $h_s\circ f_s$\\{\scriptsize $\theta_s$}};
  \node[blocco, fill=white, right=of s] (sc) {$\softmax\!\big(\cdot/\tau_s\big)$};
  \node[dato, right=of sc] (ps) {$P_s$};

  \draw[flusso] (g) -- (t);
  \draw[flusso] (l) -- (s);
  \draw[flusso] (g.east) -- ++(4mm,0) |- ([yshift=1.5mm]s.west);
  \draw[flusso] (t) -- (tc);
  \draw[flusso] (tc) -- (pt);
  \draw[flusso] (s) -- (sc);
  \draw[flusso] (sc) -- (ps);

  % loss
  \node[quant, right=10mm of $(pt)!0.5!(ps)$, align=center] (loss)
       {$\displaystyle\sum_{g}\sum_{v\neq g}\Hent\big(P_t(\vect{x}^g),P_s(\vect{x}^v)\big)$};
  \draw[flusso] (pt) -| (loss.north);
  \draw[flusso] (ps) -| (loss.south);
  \draw[grad] (loss.south) ++(-4mm,0) |- ([yshift=-3mm]ps.south) -| (s.south)
       node[pos=0.25, below, font=\scriptsize, text=SapRosso] {retropropagazione solo nello student};
  \node[font=\scriptsize, text=SapRosso, right=0.5mm of pt, yshift=3mm] {sg};

  % EMA and centering
  \draw[->, SapRossoScuro, thick] (s.north) -- (t.south)
       node[midway, right, font=\scriptsize, align=left]
       {$\theta_t\leftarrow m\,\theta_t+(1-m)\,\theta_s$};
  \draw[->, SapGrigio, densely dotted, thick] (tc.north) to[out=110, in=70, looseness=4]
       node[above, font=\scriptsize, text=SapGrigio] {$\vect{c}\leftarrow\lambda\vect{c}+(1-\lambda)\,\overline{h_t f_t(\vect{x}^{g})}$} (tc.north east);
\end{tikzpicture}
}
  \caption[Auto-distillazione DINO]{Auto-distillazione con multi-crop. Il teacher vede solo le
    viste globali e non riceve gradiente (sg). La sua uscita è centrata e resa più concentrata
    con una temperatura bassa. Il teacher è una media mobile esponenziale dello student e il
    centro è una media mobile delle uscite del teacher.}
  \label{fig:dino}
\end{figure}

\subsection{Teacher come media mobile e centering}

Il teacher non viene ottimizzato: i suoi parametri inseguono quelli dello student con una media
mobile esponenziale, alla maniera di un \emph{mean teacher}~\cite{tarvainen2017meanteacher} o di
una media di Polyak~\cite{polyak1992averaging}:
\begin{equation}
  \theta_t\leftarrow m\,\theta_t+(1-m)\,\theta_s,\qquad
  m(\tau)=1-(1-m_0)\,\frac{1+\cos(\pi\tau/T)}{2},
  \label{eq:dino-ema}
\end{equation}
con $m_0=0{,}996$ che cresce fino a 1 alla fine dell'addestramento (passo $\tau$ su $T$). Il
teacher nasce come copia esatta dello student, e la media include le statistiche della batch
normalization. Il centro è a sua volta una media mobile delle uscite del teacher sul batch:
\begin{equation}
  \vect{c}\leftarrow\lambda\,\vect{c}+(1-\lambda)\,\frac{1}{2B}\sum_{i=1}^{B}\sum_{a\in\{g_1,g_2\}}g_t(\vect{x}_i^{a}),
  \qquad\lambda=0{,}9 .
  \label{eq:dino-centro}
\end{equation}

\subsection{Collasso e sua diagnosi}
\label{sec:collasso}

L'entropia incrociata si scompone come
$\Hent(P_t,P_s)=\Hent(P_t)+\KL(P_t\,\Vert\,P_s)$. Se il teacher produce la stessa distribuzione
per ogni ingresso, lo student la può imitare perfettamente ($\KL=0$) senza aver imparato nulla:
è il \emph{collasso}. Ne esistono due forme:
\begin{itemize}
  \item \textbf{collasso su una dimensione}: $P_t$ è quasi un vettore indicatore, uguale per
  tutte le immagini, e $\Hent(P_t)\to0$;
  \item \textbf{collasso uniforme}: $P_t$ è quasi uniforme per tutte le immagini, e
  $\Hent(P_t)\to\log K$ ($\approx9{,}70$ nat per $K=\num{16384}$).
\end{itemize}
Il \emph{centering}~\eqref{eq:dino-centro} sottrae la componente comune a tutte le immagini e
impedisce a una dimensione di dominare, ma da solo spinge verso l'uniforme. Lo
\emph{sharpening} ($\tau_t$ piccola) ha l'effetto opposto. DINO bilancia le due forze.

Entrambe le forme di collasso hanno in comune l'indipendenza dell'uscita dall'ingresso, che si
misura con l'informazione mutua tra immagine e prototipo:
\begin{equation}
  I=\Hent\bigl(\E_{\vect{x}}[P_t(\vect{x})]\bigr)-\E_{\vect{x}}\bigl[\Hent(P_t(\vect{x}))\bigr]\;\ge\;0 .
  \label{eq:dino-mi}
\end{equation}
Il collasso, in entrambe le forme, corrisponde a $I\to0$. Nel monitoraggio si registra a ogni
epoca l'entropia media del teacher $\E_{\vect{x}}[\Hent(P_t)]$, che tende a $0$ o a $\log K$ in
caso di collasso.
\Segnaposto{Aggiungere al monitoraggio anche l'entropia della distribuzione media
$\Hent(\E_{\vect{x}}[P_t])$, così da stimare la~\eqref{eq:dino-mi}.}

\subsection{Schedule}

Il learning rate cresce linearmente per 10 epoche fino a $\eta_0=5\cdot10^{-4}\cdot B/256$
(regola di scala lineare) e poi decresce con un coseno fino a $10^{-6}$. Il \emph{weight decay}
cresce con un coseno da 0,04 a 0,4. La temperatura del teacher cresce linearmente da 0,04 a 0,07
nelle prime 30 epoche, mentre $\tau_s=0{,}1$ resta fissa. Uno schedule coseno da $v_0$ a
$v_{\mathrm{fin}}$ su $T$ passi è
$v(\tau)=v_{\mathrm{fin}}+(v_0-v_{\mathrm{fin}})\frac{1+\cos(\pi\tau/T)}{2}$. Il
\cref{alg:dino} riassume un passo di addestramento.

\begin{algorithm}[t]
  \caption{Un passo di addestramento DINO}
  \label{alg:dino}
  \begin{algorithmic}[1]
    \Require batch $\{\vect{x}_i\}_{i=1}^{B}$, student $\theta_s$, teacher $\theta_t$, centro $\vect{c}$, passo $\tau$
    \State $(\eta,\mathrm{wd},m,\tau_t)\gets$ schedule al passo $\tau$
    \ForEach{immagine $\vect{x}_i$}
      \State genera $\vect{x}_i^{g_1},\vect{x}_i^{g_2}$ ($224^2$) e $\vect{x}_i^{l_1},\dots,\vect{x}_i^{l_M}$ ($96^2$)
    \EndFor
    \State $P_t\gets\softmax\bigl((g_t(\vect{x}^{g})-\vect{c})/\tau_t\bigr)$ \Comment{senza gradiente}
    \State $P_s\gets\softmax\bigl(g_s(\vect{x}^{v})/\tau_s\bigr)$ per tutte le viste $v$
    \State $\mathcal{L}\gets\frac1B\sum_i\frac{1}{\abs{\mathcal{P}}}\sum_{(a,b)\in\mathcal{P}}\Hent\bigl(P_t(\vect{x}_i^{a}),P_s(\vect{x}_i^{b})\bigr)$
    \State aggiorna $\theta_s$ con AdamW($\eta$, wd) e clipping della norma del gradiente
    \State $\theta_t\gets m\,\theta_t+(1-m)\,\theta_s$
    \State $\vect{c}\gets\lambda\vect{c}+(1-\lambda)\,\mathrm{media}_{i,a}\,g_t(\vect{x}_i^{a})$
  \end{algorithmic}
\end{algorithm}

\subsection{Augmentation e immagini aeree}
\label{sec:augmentation-teoria}

Le augmentation definiscono le invarianze che la rappresentazione impara: tutto ciò che varia
tra due viste della stessa immagine viene trattato come rumore. Nelle immagini naturali le
trasformazioni cromatiche forti (variazioni di luminosità, contrasto e saturazione, scala di
grigi, solarizzazione) insegnano a ignorare l'illuminazione. Nelle immagini aeree, però, il
colore è spesso l'indizio principale: in falso colore IR-R-G la vegetazione è rossa per la
riflettanza nel vicino infrarosso. Rendere la rete invariante al colore può quindi eliminare
informazione utile. D'altra parte, le immagini aeree non hanno un'orientazione canonica, e le
rotazioni di \ang{90} e le riflessioni (il gruppo diedrale $D_4$) sono invarianze corrette che le
augmentation di DINO non usano. La variante \enquote{aerea} studiata nella \RQ{4} attenua le
trasformazioni cromatiche e aggiunge le trasformazioni di $D_4$.

\subsection{Dal pre-addestramento alla segmentazione}

Al termine del pre-addestramento si conserva solo l'encoder dello student, $f_s$, che inizializza
l'encoder di \maresunet{}. Attenzione, decoder e testa partono da un'inizializzazione casuale.

\begin{osservazione}[Obiettivo globale, compito denso]
\label{oss:globale-denso}
La perdita~\eqref{eq:dino-loss} agisce dopo la media globale $\GAP$: premia rappresentazioni
che descrivono l'immagine nel suo insieme, non le singole posizioni. Il compito a valle è invece
denso. Nulla nell'obiettivo garantisce che le mappe intermedie $\vect{E}_1,\vect{E}_2$, che
alimentano le connessioni di salto, siano informative a livello di pixel. È una possibile
ragione di un divario tra DINO e ImageNet indipendente dalla quantità di dati, e motiva la
variante densa discussa nel \cref{cap:prospettive}.
\end{osservazione}

% ---------------------------------------------------------------------
\section{Fusione dell'informazione altimetrica}
\label{sec:ndsm}

\subsection{Dal DSM all'nDSM}

Il modello digitale di superficie (DSM) dà la quota assoluta della superficie visibile, quindi
anche di tetti e chiome. Il modello normalizzato (nDSM) sottrae la quota del terreno (DTM) e
restituisce l'altezza sul terreno, che è l'informazione utile a distinguere un tetto da una
strada o un albero da un prato. Quando l'nDSM non è disponibile, il DTM si può stimare con
un'\emph{apertura morfologica}~\cite{soille2003morphology}. Detti
$\varepsilon_B(f)(u)=\min_{v\in B}f(u+v)$ ed $\delta_B(f)(u)=\max_{v\in B}f(u-v)$ l'erosione e
la dilatazione con elemento strutturante $B$, si ha
\begin{equation}
  \widehat{\mathrm{DTM}}=\gamma_B(\mathrm{DSM})=\delta_B\bigl(\varepsilon_B(\mathrm{DSM})\bigr),
  \qquad
  \mathrm{nDSM}=\mathrm{DSM}-\widehat{\mathrm{DTM}} .
  \label{eq:apertura}
\end{equation}
L'apertura elimina le strutture in rilievo più strette di $B$ (edifici, alberi, auto) e conserva
l'andamento del terreno. Essendo \emph{anti-estensiva} ($\gamma_B(f)\le f$), garantisce
$\mathrm{nDSM}\ge0$ per costruzione. In pratica si calcola il minimo del DSM su celle di circa
\SI{1}{\meter}, si applica l'apertura con una finestra di circa \SI{30}{\meter}, si smussa il
risultato e lo si riporta alla risoluzione originale. Il limite noto è che gli edifici più larghi
della finestra lasciano parte della loro altezza nel terreno stimato. L'nDSM viene codificato in
decimetri su 8 bit, $h=\min(\mathrm{round}(10\cdot\mathrm{nDSM}),255)$, con saturazione a
\SI{25,5}{\meter}.

\subsection{Fusione precoce ed estensione dei pesi}

Nella fusione precoce l'nDSM normalizzato con media e deviazione standard del training,
$\tilde h=(h-\mu_h)/\sigma_h$, diventa un quarto canale di ingresso. Il primo strato di un
encoder pre-addestrato su tre canali, con pesi $W\in\R^{k\times k\times3\times C}$, va esteso a
$W'=[\,W\mid\vect{w}_4\,]\in\R^{k\times k\times4\times C}$.

\begin{proposizione}[Estensione con pesi nulli]
\label{prop:inflazione}
Se $\vect{w}_4=\vect{0}$, la rete con ingresso $[\vect{x};\tilde h]$ calcola esattamente la
stessa funzione della rete RGB con ingresso $\vect{x}$, per ogni $\tilde h$. Il gradiente
rispetto a $\vect{w}_4$ è però in generale non nullo, quindi l'addestramento può imparare a
usare l'altezza partendo dal modello RGB.
\end{proposizione}
\begin{proof}
Per linearità della convoluzione,
$W'*[\vect{x};\tilde h]=W*\vect{x}+\vect{w}_4*\tilde h=W*\vect{x}$ quando $\vect{w}_4=\vect{0}$.
Detto $\vect{\delta}$ il gradiente della perdita rispetto all'uscita del primo strato,
$\partial\mathcal{L}/\partial\vect{w}_4[a,b,:]=\sum_u\tilde h(u+(a,b))\,\vect{\delta}(u)$, che è
non nullo appena l'altezza è correlata con l'errore.
\end{proof}

L'alternativa $\vect{w}_4=\frac13\sum_cW_{:,:,c,:}$ tratta l'altezza come un canale di luminanza
aggiuntivo. Perturba la funzione iniziale, ma dà subito al canale un peso dell'ordine di quelli
RGB. Il confronto tra le due inizializzazioni fa parte della \RQ{6}.

% ---------------------------------------------------------------------
\section{Metriche di valutazione}
\label{sec:metriche}

\subsection{Matrice di confusione e metriche per classe}

Tutte le metriche derivano dalla matrice di confusione $M\in\N^{K\times K}$, accumulata su
\emph{tutti} i pixel etichettati dello split:
$M_{ij}=\abs{\{u\in\Omega_v : y_u=i,\ \hat y_u=j\}}$. Per la classe $k$:
\begin{gather}
  \mathrm{TP}_k=M_{kk},\qquad
  \mathrm{FP}_k=\sum_{i\neq k}M_{ik},\qquad
  \mathrm{FN}_k=\sum_{j\neq k}M_{kj},\nonumber\\
  P_k=\frac{\mathrm{TP}_k}{\mathrm{TP}_k+\mathrm{FP}_k},\qquad
  R_k=\frac{\mathrm{TP}_k}{\mathrm{TP}_k+\mathrm{FN}_k},\nonumber\\
  F1_k=\frac{2\,\mathrm{TP}_k}{2\,\mathrm{TP}_k+\mathrm{FP}_k+\mathrm{FN}_k},\qquad
  \mathrm{IoU}_k=\frac{\mathrm{TP}_k}{\mathrm{TP}_k+\mathrm{FP}_k+\mathrm{FN}_k},
  \label{eq:metriche}
\end{gather}
e l'accuratezza complessiva è $\mathrm{OA}=\operatorname{tr}(M)/\sum_{ij}M_{ij}$. Precision e
recall corrispondono a \emph{correctness} e \emph{completeness} nella terminologia ISPRS.
L'IoU (indice di Jaccard~\cite{everingham2010pascal}) e l'F1 sono legati da una relazione
monotona:
\begin{equation}
  \mathrm{IoU}_k=\frac{F1_k}{2-F1_k},
\end{equation}
che si verifica sostituendo $T=\mathrm{TP}_k$ ed $E=\mathrm{FP}_k+\mathrm{FN}_k$ nelle
definizioni. Le due metriche ordinano quindi i modelli allo stesso modo classe per classe, ma
non necessariamente in media. mIoU e mF1 sono medie sulle classi presenti nel riferimento.
Vengono riportate con tutte le classi e senza clutter, che in Vaihingen è quasi assente.

\subsection{Perché una matrice di confusione globale}

Mediare le metriche calcolate batch per batch non equivale a calcolarle sull'intero split,
perché la media di rapporti è diversa dal rapporto delle somme. L'esempio della
\cref{tab:media-batch} mostra lo scostamento su una classe rara. Usare \codice{nanmean} per
saltare i batch in cui la classe è assente peggiora le cose, perché fa pesare un batch con un
solo pixel della classe quanto un batch che ne contiene migliaia.

\begin{table}[t]
  \centering
  \small
  \caption[Media per batch e matrice di confusione globale]{IoU di una classe rara calcolato
    come media per batch e dalla matrice di confusione globale.}
  \label{tab:media-batch}
  \begin{tabular}{|l|r|r|r|c|}
    \hline
    & \textbf{TP} & \textbf{FP} & \textbf{FN} & \textbf{IoU} \\
    \hline
    Batch A (un'auto piccola) & 1 & 0 & 0 & 1,000 \\
    Batch B (parcheggio) & 90 & 60 & 50 & 0,450 \\
    \hline
    Media per batch & & & & 0,725 \\
    \rowcolor{RigaRiferimento}
    Matrice globale & 91 & 60 & 50 & 0,453 \\
    \hline
  \end{tabular}
\end{table}

\subsection{Etichette piene ed \emph{eroded}}

Il benchmark ISPRS fornisce anche etichette \emph{eroded}, in cui una fascia di pixel lungo i
contorni degli oggetti è marcata come non valida. Il motivo è che la posizione esatta di un
contorno è incerta anche per l'annotatore. Valutare sulle etichette eroded significa sostituire
$\Omega_v$ con $\Omega_v^{\mathrm{er}}\subset\Omega_v$. Il confronto tra le due versioni
quantifica quanta parte dell'errore si concentra sui bordi. Le definizioni coincidono con quelle
dello strumento di valutazione ufficiale ISPRS.

% ---------------------------------------------------------------------
\section{Inferenza su tile intere}
\label{sec:finestra}

Le tile di test (fino a $6000\times6000$ pixel) non entrano in memoria in un solo passo. Si
predicono quindi con finestre di lato $w$ e passo $s<w$, e le probabilità delle finestre
sovrapposte vengono fuse. Vicino ai bordi di una finestra il contesto è troncato e la
predizione è meno affidabile: le probabilità vengono quindi pesate con una finestra di Hann
separabile,
\begin{equation}
  \omega(u)=\omega_1(u_x)\,\omega_1(u_y),\qquad
  \omega_1(t)=\sin^2\!\left(\frac{\pi(t+1)}{w+1}\right),\quad t=0,\dots,w-1,
  \label{eq:hann}
\end{equation}
con un minimo di $10^{-3}$. Dette $o_j$ le origini delle finestre e $\vect{p}_j$ le loro
probabilità,
\begin{equation}
  P(u)=\frac{\sum_j\omega(u-o_j)\,\vect{p}_j(u-o_j)}{\sum_j\omega(u-o_j)},
  \qquad
  \hat y_u=\argmax_kP_k(u).
  \label{eq:fusione}
\end{equation}
Con passo $s=w/2$ e nel limite continuo vale $\sin^2(\pi t/w)+\sin^2(\pi(t+w/2)/w)=1$: le
finestre formano una partizione dell'unità e il denominatore della~\eqref{eq:fusione} è
costante lontano dai bordi della tile (\cref{fig:hann}). La normalizzazione serve allora solo
vicino ai bordi e dove l'ultima finestra di ogni riga e colonna è allineata al bordo della tile.
Le tile più piccole della finestra vengono estese per riflessione.

\begin{figure}[t]
  \centering
  % Pesi delle finestre di Hann con sovrapposizione del 50%: la somma vale 1 lontano dai bordi.
\begin{tikzpicture}
  \begin{axis}[
      width=0.9\linewidth, height=4.8cm,
      xlabel={posizione lungo la tile (in unità di lato della finestra $w$)},
      ylabel={peso},
      xmin=0, xmax=3, ymin=0, ymax=1.25,
      xtick={0,0.5,...,3}, ytick={0,0.5,1},
      grid=major, grid style={SapGrigioChiaro},
      legend style={font=\footnotesize, draw=SapGrigioTenue, at={(0.5,1.03)}, anchor=south},
      legend columns=3,
      tick label style={font=\footnotesize}, label style={font=\small},
      samples=200,
    ]
    \addplot[SapRossoTenue, thick, domain=0:1] {sin(deg(pi*x))^2};
    \addlegendentry{$\omega$ delle singole finestre}
    \addplot[SapRossoTenue, thick, domain=0.5:1.5, forget plot] {sin(deg(pi*(x-0.5)))^2};
    \addplot[SapRossoTenue, thick, domain=1:2, forget plot] {sin(deg(pi*(x-1)))^2};
    \addplot[SapRossoTenue, thick, domain=1.5:2.5, forget plot] {sin(deg(pi*(x-1.5)))^2};
    \addplot[SapRossoTenue, thick, domain=2:3, forget plot] {sin(deg(pi*(x-2)))^2};
    \addplot[SapRosso, very thick, domain=0.5:2.5] {1};
    \addlegendentry{somma dei pesi}
    \addplot[SapGrigio, very thick, dashed, domain=0:0.5] {sin(deg(pi*x))^2};
    \addlegendentry{bordo: compensato dalla normalizzazione}
    \addplot[SapGrigio, very thick, dashed, domain=2.5:3, forget plot] {sin(deg(pi*(x-2)))^2};
  \end{axis}
\end{tikzpicture}

  \caption[Fusione delle finestre con pesi di Hann]{Pesi di Hann con sovrapposizione del 50\%:
    la somma è costante lontano dai bordi, quindi la fusione~\eqref{eq:fusione} è una media
    convessa che privilegia il centro di ogni finestra.}
  \label{fig:hann}
\end{figure}
